\documentclass[12pt]{ctexart}  % ✅ macOS 中文安全、兼容性好
\usepackage{xeCJK}
\usepackage{fontspec}
\usepackage{amsmath,amssymb,mathtools}
\usepackage{geometry}
\usepackage{hyperref}
\usepackage{enumitem}
\usepackage{titlesec}
\usepackage{microtype}
\usepackage{caption}
\usepackage{bm}

\usepackage{xcolor}
% 定义 Bronze 组的代表色
\definecolor{bronze}{HTML}{CD7F32}

% --- 代码高亮支持 (已配置好 C++ 样式) ---
\usepackage{listings}
\lstset{
    language=C++,
    basicstyle=\ttfamily\small,       % 字体大小
    keywordstyle=\color{blue}\bfseries, % 关键字颜色
    commentstyle=\color{gray},        % 注释颜色
    stringstyle=\color{red!70!black}, % 字符串颜色
    numbers=left,                     % 行号显示在左侧
    numberstyle=\tiny,                % 行号字体
    frame=single,                     % 显示边框
    breaklines=true,                  % 自动换行
    tabsize=4                         % Tab 缩进宽度
}
% -----------------------

\geometry{left=2.5cm,right=2.5cm,top=2.5cm,bottom=2.5cm}
\titleformat{\section}{\Large\bfseries}{\thesection}{1em}{}
\titleformat{\subsection}{\large\bfseries}{\thesubsection}{1em}{}

% =============================================
% 在此处修改标题中的年份和月份
% =============================================
\title{云斗 2026 年 9 月 {~\color{bronze}{Bronze Round}} 题目解析}
\date{2026 年 9 月 20 日}
\author{543tpbl}


\begin{document}
\maketitle


% =================================================================
% 示例部分：A. A+B Problem
% 请参照此格式编写后续题解
% =================================================================
\section{A. 异域旅行 (travel)}

考虑按照题意模拟。显然那个符号代表异或，你只需要对于 $n$ 个数都按照顺序求一个值即可。异或的运算符是 \texttt{$\wedge$}。


% \newpage % 强制换页，使排版更整齐（可选）

% =================================================================
% 骨架部分：请在下方填入实际题目内容
% =================================================================

\section{B. 分糖果 (candy)}

\subsection{做法一：[填入方法名，如：暴力枚举]}

显然可以直接考虑每个小朋友取的糖果数量，因为只会有一种。累加答案即可。

\subsection{测试点 1\textasciitilde 2}

这个题显然你是可以 dp 去考虑 $dp_{i,j}$ 代表前 $i$ 个点用了 $j$ 个糖果然后暴力转移的。考虑对于这个 dp 的过程去进行一些贪心。具有一些启示作用。

\subsection{测试点 7\textasciitilde 10}

首先考虑判断无解。显然 $m$ 的范围应当在 $\displaystyle \sum_{i=1}^{n} l_i \sim \sum_{i=1}^{n} r_i$ 之间。否则显然无解。我们继续考虑一下贪心的策略。

在不考虑能不能凑到 $m$ 的情况下，我们显然有贪心对于 $r_i \le \frac{m}{n}$ 的选 $r_i$，对于 $l_i \ge \frac{m}{n}$ 的选 $l_i$。这两种都不是显然选择 $\frac{m}{n}$。你发现这已经是最优的了，无论何种形式的调整都一定会增加答案，且这个增加一定是可以做到的。于是这个题就做完了。


% -----------------------------------------------------------------

\section{C.哥的拔河猜想(theory)}

题目要求对每组询问 $N,k$，在所有不超过 $N$ 的正整数中，找出能够表示成连续 $k$ 个质数之和的最大数。若不存在这样的数，则输出 $-1$。

设质数序列为 $p_1=2,p_2=3,p_3=5,\dots$。一个合法答案必须形如
$$
p_i+p_{i+1}+\cdots+p_{i+k-1}
$$
并且这个值不超过 $N$。

由于询问次数 $T<2000$，而 $N\le 10^6$，自然考虑先预处理质数序列。之后对于每组询问，只需要在连续 $k$ 个质数的和中找到不超过 $N$ 的最大值。

\subsection{做法 1：直接枚举连续段}

对于 $20\%$ 的数据，满足 $1\le N\le 100$。此时答案不会很大，可以直接筛出较小范围内的质数，然后枚举连续 $k$ 个质数的起点。

对于每个起点 $i$，暴力计算
$$
p_i+p_{i+1}+\cdots+p_{i+k-1}
$$
若这个和不超过 $N$，就用它更新答案。由于质数序列递增，连续 $k$ 个质数的和也会随起点后移而变大，所以当当前这一段已经超过 $N$ 时，后面的连续段不可能再合法，停止即可。

若第一段 $p_1+p_2+\cdots+p_k$ 就已经大于 $N$，则不存在合法答案，输出 $-1$。

设预处理出的质数个数为 $P$。每个起点需要 $O(k)$ 求和，起点数量为 $O(P)$，所以每组询问复杂度为 $O(Pk)$。期望得分 20 分。

\subsection{做法 2：前缀和优化枚举}

对于另外 $20\%$ 的数据，满足 $T=1$。此时只有一组询问，可以用前缀和把连续质数和的计算优化到 $O(1)$。

设筛出的质数序列为 $p_1,p_2,\dots,p_m$，定义前缀和 $S_i$ 为
$$
S_i=p_1+p_2+\cdots+p_i,\quad S_0=0
$$
那么从第 $i$ 个质数开始，长度为 $k$ 的连续质数和为
$$
S_{i+k-1}-S_{i-1}
$$
也可以写成以第 $r$ 个质数结尾的形式，其中 $r=i+k-1$，对应的连续段和为
$$
S_r-S_{r-k}
$$

于是对于一组询问，只需要枚举结尾位置 $r=k,k+1,\dots,m$，计算 $S_r-S_{r-k}$。若它不超过 $N$，就更新答案；若它已经超过 $N$，由于后续连续段和更大，可以停止枚举。

前缀和的作用是避免每次都重新累加 $k$ 个质数。连续段求和从原来的 $O(k)$ 降为 $O(1)$，因此枚举所有可能连续段即可。

设筛的上界为 $V=10^6$，筛出的质数个数为 $P$。筛质数复杂度为 $O(V)$，单次询问枚举复杂度为 $O(P)$，总复杂度为 $O(V+P)$。期望得分 40 分。

\subsection{做法 3：所有询问中的 $N$ 均相等}

对于另外 $20\%$ 的数据，所有询问中的 $N$ 均相等。此时不同询问之间只有 $k$ 可能不同，因此可以按照 $k$ 分组处理。

对于相同的 $k$，由于 $N$ 也相同，答案一定相同。于是只需要对每个出现过的 $k$ 计算一次答案，再把结果填回所有对应询问。

具体来说，对于某个固定的 $k$，用前缀和枚举所有长度为 $k$ 的连续质数和，即枚举
$$
S_k-S_0,\quad S_{k+1}-S_1,\quad S_{k+2}-S_2,\dots
$$
在其中找到不超过这个固定 $N$ 的最大值。如果 $k>m$，说明质数数量不足，直接无解。如果 $S_k-S_0>N$，由于这是最小的一段，后面的段只会更大，也无解。

这个做法相比对每组询问都枚举一遍，减少了相同 $k$ 的重复计算。若询问中不同的 $k$ 数量较少，就能明显加速。

设不同的 $k$ 的数量为 $Q$，筛出的质数个数为 $P$。筛质数复杂度为 $O(V)$，每个不同的 $k$ 需要 $O(P)$ 枚举，总复杂度为 $O(V+QP)$。期望得分 60 分。

\subsection{做法 4：二分查找}

对于完整数据，满足 $1\le T<2000$，$1\le N\le 10^6$，$1\le k$。如果每组询问都线性枚举所有连续段，虽然 $T$ 不算特别大，但仍然不是最稳的做法。关键在于固定 $k$ 后，所有长度为 $k$ 的连续质数和是严格递增的。

设第 $i$ 段长度为 $k$ 的连续质数和为
$$
A_i=p_i+p_{i+1}+\cdots+p_{i+k-1}
$$
那么下一段为
$$
A_{i+1}=p_{i+1}+p_{i+2}+\cdots+p_{i+k}
$$
两式相减，有
$$
A_{i+1}-A_i=p_{i+k}-p_i
$$
由于质数序列严格递增，所以 $p_{i+k}>p_i$，因此 $A_{i+1}>A_i$。这说明固定 $k$ 后，连续 $k$ 个质数和随起点后移严格递增。

用前缀和表示，以第 $r$ 个质数结尾的长度为 $k$ 的连续质数和为
$$
S_r-S_{r-k}
$$
其中 $k\le r\le m$。由于这些值严格递增，每组询问只需要二分最大的 $r$，使得
$$
S_r-S_{r-k}\le N
$$
找到这个最大的 $r$ 后，答案就是 $S_r-S_{r-k}$。

无解判断也很直接。若 $k>m$，说明预处理出的质数数量不足以组成连续 $k$ 个质数，输出 $-1$。若最小的一段 $S_k-S_0$ 已经大于 $N$，由于后面的连续段和只会更大，也输出 $-1$。

这里还需要说明筛的上界。题目中要求答案不超过 $N$，且 $N\le 10^6$。一旦某个质数本身已经大于 $10^6$，包含它的连续正数和也一定大于 $10^6$，不可能成为任何询问的答案。因此筛到 $10^6$ 已经足够。

设筛的上界为 $V=10^6$，筛出的质数个数为 $P$。预处理质数与前缀和复杂度为 $O(V)$，每组询问二分复杂度为 $O(\log P)$，总复杂度为 $O(V+T\log P)$。期望得分 100 分。


% -----------------------------------------------------------------

\section{D. 流水 (water)}

个人认为这个题很具有启示性。

\subsection{Subtask 3}

显然你只需要判断时间为 $n$ 的时候能不能流满。因为我们有 $s_i \ge 1$ 我们就能直接确定答案为 $n$。

\subsection{Subtask 1,2}

你考虑记录一个数组 $p$，$p_i$ 表示如果要把以 $i$ 为根的子树全部流满还需要的水的量，如果可以自给自足当然为 $0$。这样的话我们只需要每次从下往上跑一下 dp 就可以测试是否可以流满。如果子树的某个 $p_i$ 大于 $0$ 就加上再加 $1$，否则不管。时间复杂度可以做到 $O(n^2)$。

\subsection{Subtask 6}

考虑普遍性做法。从 Subtask 1,2 启示我们可以判断。我们发现答案具有单调性，即如果 $t$ 时间可以，$t+1$ 时间也一定可以。于是这启示我们可以直接二分答案。这样就可以做到时间复杂度 $O(n \log n)$，可以通过此题。

\end{document}
